\documentclass{article}

\usepackage{tabularx}
\usepackage{booktabs}
\usepackage{hyperref}
\usepackage{longtable}

\title{Reflection and Traceability Report on \progname}

\author{\authname}

\date{}

\input{../Comments.text}
\input{../Common.text}

\begin{document}

\maketitle

This document summarizes how feedback received during the development of
\progname{} (FitCoachAR) was addressed between Revision~0 and the Final
Documentation, reflects on the design iterations, and reviews how the project
unfolded relative to its original plan. Each feedback item below is linked to
the closed GitHub issue that tracked it on the project repository
(\href{https://github.com/mmdmirh/cas741}{mmdmirh/cas741}).

\section{Changes in Response to Feedback}

All feedback was received either through GitHub issues or as annotated PDFs
attached to review-request issues. Every item was addressed and every
corresponding issue was closed. The tables below record the source, the
summary of the feedback, the response, and a pointer to the issue.

\subsection{Problem Statement}

\begin{longtable}{p{0.06\textwidth} p{0.18\textwidth} p{0.34\textwidth} p{0.32\textwidth}}
\toprule
\textbf{\#} & \textbf{Source} & \textbf{Feedback} & \textbf{Response}\\
\midrule
\endhead
1 & Dr.~Smith (\href{https://github.com/mmdmirh/cas741/issues/1}{\#1}) &
The reader needs to know what FitCoachAR is and what problem it solves. &
Rewrote the opening paragraph to describe the system and the gap it addresses
(affordable, hardware-free AR form feedback).\\
1 & Dr.~Smith (\href{https://github.com/mmdmirh/cas741/issues/1}{\#1}) &
CAS~741 is not a real-time systems course; drop the real-time framing. &
Removed all language framing the system as a real-time system; kept only the
scientific-computing aspects (pose estimation, angle calculation, validity
classification).\\
1 & Dr.~Smith (\href{https://github.com/mmdmirh/cas741/issues/1}{\#1}) &
Opening LaTeX quotes are rendered incorrectly (should use \texttt{``}). &
Corrected all opening quotes in the Problem Statement.\\
1 & Dr.~Smith (\href{https://github.com/mmdmirh/cas741/issues/1}{\#1}) &
A machine-learning report would be a good fit as an extra. &
Updated \texttt{Repos.csv} to list the Machine Learning Report as one of the
two extras for the project.\\
\bottomrule
\end{longtable}

\subsection{SRS}

\begin{longtable}{p{0.06\textwidth} p{0.18\textwidth} p{0.34\textwidth} p{0.32\textwidth}}
\toprule
\textbf{\#} & \textbf{Source} & \textbf{Feedback} & \textbf{Response}\\
\midrule
\endhead
5 & Y.~Mei (\href{https://github.com/mmdmirh/cas741/issues/5}{\#5}) &
Symbols \$, $\theta_{thresh}$, $\theta_{low}$ missing from the Table of
Symbols. &
Added the missing symbols.\\
6 & Y.~Mei (\href{https://github.com/mmdmirh/cas741/issues/6}{\#6}) &
DD1 and DD2 cite different sources for the same underlying fact. &
Unified the source of both DDs to \emph{Standard Biomechanics}.\\
7 & Y.~Mei (\href{https://github.com/mmdmirh/cas741/issues/7}{\#7}) &
R1 is not visibly traced to the Camera assumption in Table~R\_trace. &
Updated the SRS Traceability Matrix so R1 is explicitly linked to A\_Camera
(A2).\\
8 & Y.~Mei (\href{https://github.com/mmdmirh/cas741/issues/8}{\#8}) &
Add explicit vector-subtraction formulas for the joint-angle DDs. &
Added the vector forms
($\vec{v}_{femur} = P_{hip} - P_{knee}$, etc.) to GD\_KneeAngle and
GD\_ElbowAngle.\\
9 & Y.~Mei (\href{https://github.com/mmdmirh/cas741/issues/9}{\#9}) &
``Not occluded'' in A\_Visibility is ambiguous. &
Rewrote A\_Visibility to require a \emph{direct line of sight} to the target
joints, and propagated this phrasing to GD\_ElbowAngle during the final pass.\\
10 & Y.~Mei (\href{https://github.com/mmdmirh/cas741/issues/10}{\#10}) &
Missing ``Data Definitions'' header; leftover template Appendix; misaligned GD
tables; inconsistent TM numbering. &
Added the header, removed the template Appendix, switched GD tables to
\texttt{tabularx} for proper wrapping, and fixed the TM references.\\
4 & Dr.~Smith (\href{https://github.com/mmdmirh/cas741/issues/4}{\#4}) &
Detailed comments on the annotated SRS PDF. &
Applied throughout the revision; individual items are tracked in the commits
cited by the issue resolutions above.\\
\bottomrule
\end{longtable}

\subsection{VnV Plan}

\begin{longtable}{p{0.06\textwidth} p{0.18\textwidth} p{0.34\textwidth} p{0.32\textwidth}}
\toprule
\textbf{\#} & \textbf{Source} & \textbf{Feedback} & \textbf{Response}\\
\midrule
\endhead
13 & Y.~Mei (\href{https://github.com/mmdmirh/cas741/issues/13}{\#13}) &
Formatting inconsistencies. & Fixed.\\
14 & Y.~Mei (\href{https://github.com/mmdmirh/cas741/issues/14}{\#14}) &
Section~3.1 should list the actual reviewer names, not just roles. &
Section~3.1 now names the assigned reviewers explicitly.\\
15 & Y.~Mei (\href{https://github.com/mmdmirh/cas741/issues/15}{\#15}) &
Dataset description is too vague. &
Expanded the Dataset section to state that it covers a diverse distribution
of correct/incorrect form, varying user heights, and camera angles.\\
16 & Y.~Mei (\href{https://github.com/mmdmirh/cas741/issues/16}{\#16}) &
Traceability matrix is missing. &
Added the requirements--tests traceability matrix. It was further corrected
during the final pass so that R1--R5 match the R\_Input/R\_Process/R\_Calc/
R\_Verify/R\_Output identifiers used in the SRS.\\
17 & Y.~Mei (\href{https://github.com/mmdmirh/cas741/issues/17}{\#17}) &
Mentioning code reviews under \emph{Unit Testing} is confusing; it is already
covered under implementation verification. &
Removed the code-review discussion from the unit testing section.\\
12 & Dr.~Smith (\href{https://github.com/mmdmirh/cas741/issues/12}{\#12}) &
Detailed comments on the annotated VnV Plan PDF. &
Incorporated throughout the revision, including the restructuring noted
above.\\
\bottomrule
\end{longtable}

\subsection{MG and MIS}

\begin{longtable}{p{0.06\textwidth} p{0.18\textwidth} p{0.34\textwidth} p{0.32\textwidth}}
\toprule
\textbf{\#} & \textbf{Source} & \textbf{Feedback} & \textbf{Response}\\
\midrule
\endhead
19 & Y.~Mei (\href{https://github.com/mmdmirh/cas741/issues/19}{\#19}) &
MG is missing a revision history. &
Added a full revision-history table (v1.0--v1.3) at the top of the MG.\\
20 & Y.~Mei (\href{https://github.com/mmdmirh/cas741/issues/20}{\#20}) &
Empty list of figures; unusual section numbering (7.0.1 style). &
Added real figures (architecture diagram, TikZ Uses-Hierarchy DAG) and fixed
the heading hierarchy so numbering is well-formed.\\
21 & Y.~Mei (\href{https://github.com/mmdmirh/cas741/issues/21}{\#21}) &
M4, M5, M8 should not be classified as \emph{Software Decision}; they describe
externally visible behaviour. &
Partially applied. M4 (Pose Tracking) and M5 (Kinematic Engine) were
reclassified as \emph{Behaviour-Hiding} as suggested. M8 (Signal Smoothing)
was kept under \emph{Software Decision} because the filtering algorithm
(Kalman filter, see AC4) is not in the SRS and is a software-level design
choice. Section~6 of the MG was restructured so the subsection headers
(Hardware-Hiding / Behaviour-Hiding / Software Decision) match the table.\\
22 & Y.~Mei (\href{https://github.com/mmdmirh/cas741/issues/22}{\#22}) &
Inconsistency in M8 usage: the DAG and the MIS disagreed on the direction of
the M4/M8 dependency. &
Corrected the DAG so it shows M4 \emph{uses} M8; updated the M8 MIS so its
\emph{Uses} field reads \texttt{None} (M8 is a stateless filter).\\
23 & Y.~Mei (\href{https://github.com/mmdmirh/cas741/issues/23}{\#23}) &
M1, M2, M7 had no MIS sections. &
Added full MIS sections for M1 (\texttt{get\_frame}), M2 (\texttt{draw\_overlay},
\texttt{display}) and M7 (\texttt{update\_ui}, \texttt{get\_json}), each with
module identifier, uses, syntax and semantics.\\
24 & Y.~Mei (\href{https://github.com/mmdmirh/cas741/issues/24}{\#24}) &
Input frames are mis-listed as environment variables for M3 and M4. &
Reclassified frame inputs as access-program \emph{parameters}
(\texttt{frame: bytes}). Environment variables are now reserved for truly
external, persistent state (webcam hardware, browser DOM).\\
25 & Dr.~Smith (\href{https://github.com/mmdmirh/cas741/issues/25}{\#25}) &
Detailed comments on the annotated MG and MIS PDFs. &
Applied throughout the revision; residual issues surfaced during the final
pass (requirement-to-module trace table, M8 Uses) were also fixed.\\
\bottomrule
\end{longtable}

\subsection{VnV Report}

\begin{longtable}{p{0.06\textwidth} p{0.18\textwidth} p{0.34\textwidth} p{0.32\textwidth}}
\toprule
\textbf{\#} & \textbf{Source} & \textbf{Feedback} & \textbf{Response}\\
\midrule
\endhead
26 & Y.~Mei (\href{https://github.com/mmdmirh/cas741/issues/26}{\#26}) &
Peer reviewer was unable to schedule a VnV Report review in time. &
Acknowledged; the report will be reviewed later by the instructor
(see~\#27).\\
27 & Dr.~Smith (\href{https://github.com/mmdmirh/cas741/issues/27}{\#27}) &
Will review the VnV Report when the full project is complete. &
The report reflects the executed test plan, including the consolidation of
T3 into T2 and the decision not to implement the Robot-Framework E2E suite
originally labelled~T9.\\
\bottomrule
\end{longtable}

\section{Extras}

Two extras were pursued for this project: (i)~a \textbf{Machine Learning
Report}, replacing an earlier extra at Dr.~Smith's suggestion (issue~\#1), and
(ii)~a \textbf{User Manual} describing how to install, calibrate, and use the
system. Both extras were approved as part of the Problem Statement and are
listed in \texttt{Repos.csv}.

\section{Design Iteration (LO11)}

The final design is the result of three substantial iterations.

\textbf{Rev~0.} The first architecture treated FitCoachAR as a monolithic
browser-side pipeline in which MediaPipe, the kinematic computation, and the
state machine all lived in the same frontend process. This was convenient
for the Proof-of-Concept but made the scientific logic (angle calculation,
form validity) hard to test in isolation.

\textbf{Rev~1 (MG/MIS presentation).} After the modular-design lectures and
Dr.~Smith's feedback, the system was decomposed into eight modules along
information-hiding lines, splitting the browser (M1, M2, M7) from the
Python backend (M3--M6, M8). This made it possible to unit-test M5
(kinematic engine) and M6 (state machine) in pytest without a camera in the
loop, and directly addressed the review comment that the scientific
computing \emph{secrets} were not being surfaced.

\textbf{Rev~1.x (Final pass).} During peer review Yibing pointed out that
the ``behavioural'' modules (M4, M5, M8) had been mis-filed under
\emph{Software Decision}. The hierarchy was corrected to reflect the fact
that pose tracking, angle computation, and smoothing are all required
externally visible behaviours, not hidden implementation choices. A late
internal audit also caught a stale requirement-to-module traceability table
in the MG, which was re-derived from the SRS R\_Input\ldots R\_Output
definitions. These changes did not alter the code; they aligned the
documentation with what the system actually does.

In all three iterations the same underlying need drove the design: make
the scientific-computing core testable and auditable, rather than hiding it
behind the UI.

\section{Design Decisions (LO12)}

\textbf{MediaPipe as the pose backend (M4).} MediaPipe was chosen because it
runs on commodity webcams without depth hardware, satisfying the
``affordable, no specialized hardware'' constraint from the Problem
Statement. The decision is recorded as anticipated change \texttt{acMLModel}
in the MG so that a future replacement (e.g.\ OpenPose, BlazePose-GHUM)
would only touch M4.

\textbf{Explicit Kalman smoothing as a separate module (M8).} MediaPipe's
per-frame landmarks are noisy enough that the raw signal produces unstable
angle estimates and spurious state-machine transitions. A stateless
smoothing filter was introduced as its own module so the choice of filter
(currently a 1-D Kalman per coordinate) can change without affecting the
kinematic engine. Its \emph{Uses} set is \texttt{None} for the same reason.

\textbf{State-machine based rep counting (M6).} A threshold-hysteresis FSM
was preferred over an ML classifier because the squat and bicep-curl rep
definitions are already expressible in closed form, and an FSM gives
reviewers a clearly verifiable specification (see~IM\_RepetitionCount).
Assumption A\_Visibility and the ``direct line of sight'' constraint
followed directly from this decision: the FSM has no recovery path for
missing landmarks, so the assumption is an explicit limitation, not an
accident.

\textbf{Separation of the Python backend from the browser front-end.} This
was influenced by the testability constraint. By pushing M3--M6 and M8 into
a Python process, the scientific computing modules can be exercised by
pytest fixtures without any camera or browser, which is what T1, T2, T4--T8
actually do.

Limitations that shaped these decisions were: the inability to assume any
depth sensor, the absence of a training dataset large enough to train a
classifier of our own, and the course's emphasis on scientific-computing
verifiability over end-to-end system testing.

\section{Reflection on Project Management (LO24)}

\subsection{How Project Management Compared to the Development Plan}

The project largely followed the planned workflow: weekly self-imposed
milestones aligned with the CAS~741 deliverable schedule (Problem Statement
$\rightarrow$ SRS $\rightarrow$ VnV Plan $\rightarrow$ MG/MIS $\rightarrow$
Implementation $\rightarrow$ VnV Report $\rightarrow$ Final
Documentation). GitHub issues were used both for peer review and for
tracking Dr.~Smith's feedback, and each issue was closed with a commit
reference, which made the final traceability exercise (this document)
substantially easier.

The technology choices (Python~+ FastAPI backend, MediaPipe for pose,
browser-side three.js front-end, LaTeX documentation, pytest) stayed
consistent from the Problem Statement onward.

\subsection{What Went Well}

\begin{itemize}
  \item \textbf{Issue-driven feedback.} Every review comment from Yibing and
  Dr.~Smith became a GitHub issue, which was closed by a commit that named
  the change. This made the Revision~0 $\rightarrow$ Final Documentation
  diff trivial to reconstruct.
  \item \textbf{Early modular split.} Moving the scientific core into its
  own Python backend paid off repeatedly: it made it possible to unit-test
  the kinematic engine and state machine without a camera.
  \item \textbf{Reviewing as a source of design improvements.} Several
  issues (\#21, \#22, \#24) were not just formatting fixes; they exposed
  real inconsistencies between my mental model and what was documented. The
  reviews were the single biggest source of design-level improvements.
\end{itemize}

\subsection{What Went Wrong}

\begin{itemize}
  \item \textbf{Documentation drift.} Between Rev~0 and the final pass,
  several documents (most visibly the MG traceability table and the VnV
  Plan requirement numbering) drifted out of sync with the SRS. The
  drift was only caught by a deliberate end-of-project audit; the regular
  review cycles had not checked cross-document consistency.
  \item \textbf{Late unit testing.} Unit tests were written close to the
  implementation deadline rather than in parallel with each module. This
  meant bugs surfaced at the end when there was little time to redesign.
  \item \textbf{VnV Report scheduling.} The VnV Report was finished too
  close to the peer-review window, which made it impossible for the peer
  reviewer to return feedback in time (issue~\#26).
\end{itemize}

\subsection{What Would I Do Differently}

\begin{itemize}
  \item Run a \emph{cross-document consistency check} at every revision
  boundary, not just at the end. A short checklist (requirement IDs,
  module IDs, test IDs, assumption labels) would have caught the late
  drift much earlier.
  \item Write the unit tests alongside each module, not after. This would
  have made the VnV Report's code-coverage section much stronger.
  \item Budget the VnV Report the same amount of calendar time as the SRS
  or MG, so that peer review is actually feasible.
\end{itemize}

\end{document}
